\section{Conclusion of the Prototyping Process}

As the goal of this chapter was not just to create a usable artifact for its own right
and explore the process of building a kubernetes-based \ac{HPC} tool,
but also to stand as the validation step of the system designed in the \Nameref{suggested_architecture}.

While each iteration of the prototyping process is split into goal, implementation, and evaluation
sections, this section serves to do this for the process as a whole and conclude not just the prototyping process,
but the conceptual design of the system as a whole,
by reevaluating the original design with the insights gained from the prototyping process.


\subsection{Summary of the Prototyping Process}

Before the three iterations of the prototyping process are set up, executed and evaluated,
the existing infrastructure, previous work, and the environment of the prototyping process are described.
This is done to set the stage for the prototyping process, as well as to
highlight the first differences between the suggested architecture and the prototype to be.
This is important as the architecture is designed from a deductive standpoint,
being a system that is supposed to be generic and not influenced by previous decisions.
Confronting this design with the first practical steps is the first step in the validation process.
From this the changes extensions and removals of the previous system are derived.

The fist iteration (\ref{sec:iteration_no1}) of the prototyping process is to create the most basic version of the system,
possible, to test the feasibility of the system and the environment.
This is done by verifying the operative state of the system and the environment,
fixing issues with the existing infrastructure.
For the basic functionality of the system, a custom container image is created,
containing the necessary tools and libraries to host a Notebook server.
Using a modified \textit{Ipython} kernel, the Arkouda and Chapel libraries are
integrated into the container which automatically connect to the compute nodes when
activated. To make the system able to be used on the cluster a deployment and a
service are created, to make the system accessible from the outside.
The system is then tested by running a simple example, to verify the functionality of the system,
some issues are then identified and workarounds are implemented to create the first working version of the system.

As the first iteration  has successfully implemented a \ac{MVP} of the system,
second iteration (\ref{sec:iteration_no2}) aims to improve the system as
the previous version did not regard speed and stability as a priority.
During this iteration the previously created notebook container is rebuilt in
order to optimize the performance of the system.
The infrastructure from the previous iteration is then given the same treatment,
completely rebuilding the entire container images from scratch in order to
reduce their size and to remove version mismatches in the underlying libraries.
To measure these improvements the deployment times of the system are benchmarked.
Additionally it was attempted to replace the \ac{CRI} with a faster \ac{HPC} oriented
one, which was not successful due to kubernetes dropping support for the \ac{API}
specification which the \ac{CRI} was based on.
Even though the \ac{CRI} could not be replaced, the system was still benchmarked
to measure the compute node's performance and to identify bottlenecks in the system
For this purpose a simple benchmarking framework was converted to a deployment
capable of running in a kubernetes cluster and to interact with the compute nodes
to measure the performance of the system.

The goal of the third iteration (\ref{sec:iteration_no3}) was to move the system from
a simple proof of concept to a system which could not just technically solve \ac{HPC} problems,
but one that is not constrained by consumer level hardware, so its performance can be
compared to an equivalent section of a real system.
While the previous section had the goal of implementing \textit{Singularity} to
make use of its near native performance, it was also hoped to make use of its
native support for \ac{IB} to make the third iteration much easier to implement.
But as \textit{Singularity} was not supported the \ac{IB} was implemented manually
using a multitude of different tools and configurations in order to properly provision
and integrate the device into the Multi \ac{CNI} setup, installed to enable the usage.
During the integration of this step the prototyping was cut short due to the time
constraints of the thesis, leaving the \ac{IB} integration for future work.
But as the flannel \ac{CNI} is still completely in place the system is still
fully functional and can be used to run \ac{HPC} workloads using the \ac{TCP/IP} protocol.


\subsection{Evaluating against the Deductively Designed System}

As the prototyping process is now concluded and summarized,
it is time to complete the validation of the deductively designed system
and to evaluate the prototyping process against the original design;
contrasting the theory with practical implementations, which changes
where made out of incorrect assumptions and which where made out of necessity,
and which are improvements of the general design.

The following section will go through the different aspects defined in the
\Nameref{suggested_architecture} and evaluate them.


\textbf{\hyperref[asp:client]{Client}}:

In the design it was assumed that the client of the user would
be a local interface of Jupyter notebook, which would be connected to the
server via a web interface.
This assumption has been implemented int the prototype and even used during the
data analysis of the prototyping process.
Beyond having a local interface, such as a local Jupyter notebook or the \textit{VSCode} extension,
the system can just as easily be accessed from a web browser,
making it highly accessible and and flexible in its use.
The original design has been validated by the prototyping process.

\textbf{\hyperref[asp:language]{Language}}:

In the original design it was decided to make this project Arkouda first,
as it is a modern and easy to use way of interacting with the \ac{HPC} system.

While Arkouda itself is very user friendly and easy to use, the installation,
needing to have a precompiled Chapel version, and then manually compiling
and importing the Arkouda library, is less so.

While this made the creation of custom containers more complicated,
form the perspective of the user, the system is now as easy to use as
native python on any jupyter notebook server, making the system very user friendly.

From the perspective of compute performance, the system was as efficient as
the underlying infrastructure permitted, being bottlenecked by the
networking (see \Nameref{sec:performance_testing}).

There is no reason to discount the original decision to use Arkouda as the
primary language of the system.


\textbf{\hyperref[asp:cluster]{Cluster}}:

The original design assumed that the system should be run on a Kubernetes cluster,
while this is an integral part of the ability to scale the system,
the prototyping has shown that Kubernetes is not build for connecting
easily to hardware devices which are standard in \ac{HPC} environments.

Integrating these devices is a major challenge and when deciding on
building a cloud native \ac{HPC} system, it should be strongly considered
wether a Kubernetes cluster is the right choice for the system,
or whether the other opportunity discussed in this section should be used.
This being the bare the orchestration of bare metal systems.

So while the design inherently relies on the use of a Kubernetes cluster,
and therefore it needs to be included in the system, this is a
decision that should only be made if the high flexibility
and scalability of containers is needed.

\textbf{\hyperref[asp:jupyter_server]{Jupyter Server}}:

The original design discuses the decision between using a simple Jupyter server
or a whole JupyterHub setup, the prototyping process has shown that
the setup of a single Jupyter server is sufficient for the system,
and quite uncomplicated to set up.

Orchestrating a single or a few JupyterHub in oder to manage multiple users
would increase the complexity of the system, while the use of a single Jupyter server
per user is a much simpler solution, which is sufficient for the current goals of the system.

This decision is validated by the prototyping process.

\clearpage

\textbf{\hyperref[asp:compute_services]{Compute Services}}:

As Arkouda was already selected the choice of compute service was obvious,
as the Arkouda Client is build to interface with a Chapel backend, through the
Arkouda server witch wraps the Chapel library.

While the same limitations exist for the Arkouda server as for the
Arkouda client, needing to compile Chapel and Arkouda manually, making
an iterative development process more complicated, the experience during
runtime is user-friendly and and relatively quick, considering the
networking limitations of the system.

The original design has been implemented as described and is a valid choice for the system.
But further optimization of the system could be considered to ease and speed up the
development process, as many steps of the initial setup are still required to be done manually.
A more automated and consolidated setup would be beneficial for both the scalability and the
maintainability of the system.

Over all, the original design is validated by the prototyping process, and further potential
has been identified.

\textbf{\hyperref[asp:resource_provisioning]{Resource Provisioning / Scaling}}:

Original goal of the system in the design was to dynamically scale the system,
either on a job by job basis or by scaling the based on sessions.

While the system can be rescaled via a single variable in between sessions,
this feature has not yet been implemented deeply enough in the system.
It's implementation during later iterations was a significant factor in the decisions
made during the earlier stages of the prototyping process, such as setting
up the self configuring implementation of \ac{IB}, the optimization of the container images,
and the one notebook per user setup.

While this feature has been deemed important for the system the decision to
make all parts scalable from the beginning should be reconsidered.
Creating a system that is at first not scalable, but completely functioning
with all the relevant features and then optimizing the necessary parts for
scalability. Would have resulted in more insights early on in the prototyping process,
but would have created more work over all.
Under the light of the time constraints of the thesis, the decision to
prioritize the scalability in each step component of the system has lead to
not reaching the goal of a fully scalable system, in the time frame of the thesis.

Therefore, while the original design is valid, it should have been differently prioritized.


\textbf{\hyperref[asp:operating_system]{Operating System}}:

The original design assumed that the system would be run on a Linux based system,
which is a necessary and valid assumption, for both Kubernetes and \ac{HPC} systems.
This part is validated by the prototyping process.

A second consideration was wether to use a generic Linux distribution or a specialized one,
while the choice of a generic distribution has been sufficient for the prototyping process,
choosing a licensed distribution specialized in this sector of Kubernetes would have
enabled the access to more specialized support and documentation, which would have
speed up the prototyping process significantly.
On the other hand the use of a non generic distribution entails the reliance on
licensed software, which would have made the system less accessible to the general public.
And setting up the entire cluster from a blank system offers a much deeper understanding
of the system as a whole, which is invaluable for the development of the system.

Therefore both decisions are equal valid and the choice between them should be made
depending on the goals of the developer, for a more flexible and public system a generic
distribution would be the better choice, while for a licensed and specialized system
will provide benefits for production systems.

This decision is validated by the prototyping process.


\textbf{\hyperref[asp:container_runtime]{Container Runtime}}:

During the design phase three different container runtimes where picked out as
possible candidates on which the system could be run.

On the baseline either \ac{CRI-O} or \textit{Containerd} where considered with
\ac{CRI-O} being the preferred choice due to its closer accordance with the
\ac{OCI} specification.
As an additional option \textit{Singularity} was considered as a runtime
for the compute pods to use its performance benefits.

As the previous system was already running on \textit{Containerd} and the
benefits of \ac{CRI-O} where not significant enough to warrant the change,
the decision was made to keep the system on \textit{Containerd}.
Not invalidating the original proposal but diverting from it out
of practicality.

The decision to include \textit{Singularity} as a secondary runtime,
on the other hand was inherently misinformed, as is described in the
section \Nameref{sec:container_runtime}.
Going down this path is not recommended, as either an unsupported version
of Kubernetes would have to be used, or the \textit{Singularity} \ac{CRI}
would need to be patched.

In conclusion the original design is validated in regards to the base container runtime,
but the decision to include \textit{Singularity} has been completely invalidated.

\textbf{\hyperref[sec:container_image]{Container Image}}:

As touched upon in the compute services section, the container images used where
based on the ones provided in the contributions repository of the Arkouda project.

While they have been an excellent starting point, the prototyping process has shown that
the images where able to be be significantly reduced in size and optimized for the system,
which has improved the performance of the system.

One significant change needed to be made in the future is to create a set of \ac{IB}
capable images, as this needs to be set up during compilation time and can not be
changed during runtime.

The original design has been validated by the prototyping process and has been improved upon,
and still has potential for further improvements.

\textbf{\hyperref[asp:network_abstraction]{Network Abstraction}}:

During the original design the tradeoff common network abstraction layers make
between performance and flexibility was highlighted and it was concluded that a
multi \ac{CNI} setup would be an essential part of the system.

Creating two separate network, one for the orchestrating communication, and one
to provide the high-speed \ac{RDMA} capable communication the \ac{HPC} components
require.

It was already recognized that for this setup to work a meta \ac{CNI} such
as \textit{Multus} would be necessary to enable the use of both networks.

While the need for this system has been thoroughly shown during the prototyping process,
and the implementation of a dual \ac{CNI} setup was has been proven as the only valid solution;
Other solutions such as keeping the original \textit{Flannel} \ac{CNI} and
facilitating the \ac{RDMA} communication on the container runtime level where considered but
have been proven to be infeasible.

While the original proposal on the network abstraction level has been implemented, via
the use of a dual \ac{CNI} setup, the ramifications for the configurations of the underlying
hardware where severely underestimated, which will be discussed in more detail in the paragraph
\hyperref[par:network_interconnect]{"Network Interconnect"} of this section.

Overall using a dual \ac{CNI} setup is a valid solution if the requirements of the system
demand a high-speed \ac{RDMA} network, a containerized environment and a flexible communication
infrastructure.

If this is not the case, a bare metal ethernet network, running the workload directly on the system
or a fixed network configuration should be considered instead.

\textbf{\hyperref[asp:compute_node_hardware]{Compute Node Hardware}}:

In the effort to create a generic design, only generic recommendations where made
during the design phase, especially as the system, using containerization is
largely hardware agnostic.

As the network revealed itself to be the main bottleneck of the system, the
compute and memory specification of the system did not impact the performance
in a significant negative way.
This is demonstrated by the fact that during the performance testing \ac{PGAS}
memory occupation did not exceed 30\% of the available space when
the 10Gbit network was at its maximum capacity.

But nevertheless, the memory speed and capacity, as well as the \ac{CPU} performance and \ac{NUMA}
topology of the two socket nodes would most likely be of significant importance
after the implement of a dedicated \ac{IB} network.

The other aspect discussed, being the capability of respecting non homogeneous
systems was important during the building and testing of the network abstraction.
Using a hardware scanning daemonsets, to identify available hardware made it
easy to segment the deployments to either the \ac{IB} capable or in capable nodes,
by simply stating the need in the configuration.

While not proven during the prototyping process, the assumption that a
more compute capable system would be able to handle more workloads still stands.
On the other hand it was shown that the use of heterogeneous systems is possible,
as special need hardware can be reserved for special workloads if needed,
but is fully capable of running pods with on all nodes, independent of
the underlying hardware.

Validating the latter and not invalidating the former, part of the original design.

\textbf{\hyperref[sec:suggarch_network]{Network Interconnect}}:
\label{par:network_interconnect}

As has previously been referenced, one of the major bottlenecks of the prototype was
the network. This was already correctly anticipated during the design phase and
the decision to utilize an \ac{IB} network interconnect was made.

During the design phase it was assumed that the installation of the \ac{IB} network,
would be of equal complexity as any other \ac{CNI} addon; not trivial but manageable.

But this notion was shown to be incorrect, as the setup of the \ac{IB} network
and the necessary multi layered infrastructure to configure it through Kubernetes
was a major challenge took up a significant amount of time during the prototyping process.

While the Kubernetes side of the setup has now been brought to a working state,
the configuration of the necessary kernel modules has been cut short due to
time constraints of the thesis.

While the reason for the \ac{IB} interconnect is valid; bringing the performance from
acceptable for medium sized workloads, to good performance for large scale \ac{HPC} workloads.
The complexity of this setup was severely misjudged and resulted in a major disruption
of the time schedule of the prototyping process, resulting in features not being implemented
and the original design not being able to be fully validated.

One of these features is the dynamic scaling of the system, which would have lead to
the investigation of the need for topological awareness, which has also been discussed
in this section of the design.

While the need for a high-speed network interconnect has manifested as expected,
the applied solution was not feasible in the time frame of the thesis.
Instead it is recommended to either circumvent the need for a \ac{IB} integration,
or to allot this step the necessary resources to be implemented correctly.
This could be done by making use of enterprise support, hiring a specialist
or allotting more time to the prototyping process.
Regarding the network topology, no evaluation can be made, as neither
a state in which the need could be evaluated nor a test setup has been reached.


\textbf{Conclusion}

Circling back to the original goal of the Conceptual Deductive Design,
as defined in \Nameref{sec:goals}:

\textit{\say{The designed system should provide the user with a easy to use interface,
        capable of interactively exploring HPC scale data.
        The system should be able to run on cloud infrastructure, be able to
        to allocate the necessary resources, as well as scheduling and executing
        the computation within a reasonable time frame.}}

This prototyping process has shown that this goal is technically feasible,
and has largely been achieved. While two of the design goals where based on
incorrect assumptions, most have either been validated or even expanded
upon during the prototyping process.

These divergences between a theoretical design and the realities of an
actual implementation, have now been either been corrected, put in
a more appropriate context or have been proven to be based on incorrect
assumptions.

Overall the design as well as the prototyping process has been a success,
as the central functionality of an interactive system was successfully implemented,
but most importantly valuable insights have been gained about the prototyping process
as well as the cavitates of a deductively designed system.
Also a useful system has been created which was used to run the data analysis
While there is still work to be done and much long term potential for
the project as a whole, what has been done builds a solid foundation
and many paths forward.
